\section{Formal distributions}
    \subsection{Calculus with formal distributions}
        \begin{definition}[Formal distributions] \label{def: formal_distributions}
            For any $k$-algebra $A$, which needs not be commutative, let us write:
                $$A[\![z_1^{\pm}, ..., z_n^{\pm}]\!] := \bigotimes_{i \in \{1, ..., n\}} A[\![z_i^{\pm 1}]\!] \cong \bigotimes_{i \in \{1, ..., n\}} ( A[\![z_i]\!] \oplus z_i^{-1} A[\![z_i^{-1}]\!] )$$
            for the vector space of \textbf{formal distributions} in the variables $z_1^{\pm 1}, ..., z_n^{\pm 1}$ with coefficients in $A$.
        \end{definition}
        \begin{remark}
            We caution the reader that \textit{a priori}, $A[\![z_1^{\pm}, ..., z_n^{\pm}]\!]$ is not even a $k$-algebra, since the coefficients of the product of two arbitrary elements can be infinite products of the coefficients of the two factors.  
        \end{remark}
        \begin{example}[Delta-distributions] \label{example: delta_distributions}
            A very useful class of formal distributions are the so-called \textbf{delta-distributions}:
                $$\1(z, w) := \sum_{m \in \Z} z^m w^{-m - 1}$$
            \end{example}

        \todo[inline]{Technical subsection. Add more details as the other sections progress.}

    \subsection{Quantum fields and locality}
        For what follows, let us recall the following fundamental idea from physics: in quantum field theory, particles arise as \say{excitements} of so-called \say{quantum fields}, and the motions of such particles throughout spacetimes are governed also by the nature of those quantum fields. Should one equate particles with their wave functions - which themselves can be thought of as vectors - then the aforementioned process of exciting the quantum fields can then be thought of as actions on wave functions. Suppose also, that particles exist pointwise in some sense. Quantum fields themselves, hence, can be thought of as families of operators of \say{quantum observables}, varying over points of a chosen spacetime.

        \begin{convention}
            For us, any possible spacetime will locally around points resemble the formal disc $\Spec k[\![z]\!]$, whose function field\footnote{Or field of globally meromorphic functions, if you are a complex geometer/analyst.} is $k(\!(z)\!)$. Definition \ref{def: quantum_fields}, in particular, is only reasonable under this assumption.
        \end{convention}
    
        \begin{definition}[Quantum fields] \label{def: quantum_fields}
            A \textbf{quantum field} on a given vector space $V$ (to be thought of as the space of quantum states) is a Laurent series in one variable with coefficients in $V$, i.e. an element of:
                $$\quantumfields(V) := \End_k(V)(\!(z)\!)$$

            A given quantum field:
                $$A(z) := \sum_{i \in \Z} A_i z^{-i} \in \quantumfields(V)$$
            is said to be \textbf{homogeneous} and of \textbf{conformal dimension} $d \in \Z$ if each coefficient $A_i \in \End_k(V)$ is a homogeneous operator of degree $-i + d$. As a shorthand, we will be writing:
                $$\dim A(z) = d$$
            in such a situation, pre-supposing that $A(z)$ is homogeneous.
        \end{definition}
        \begin{remark}[Formal derivative raise conformal dimensions]
            Consider a homogeneous quantum field $A(z) := \sum_{i \in \Z} A_i z^{-i}$. One notices easily that if $\dim A(z) = d$ then:
                $$\dim \del_z A(z) = d + 1$$
        \end{remark}

        \begin{question}
            Let $V$ be a vector space (of quantum states) and let $A(z), B(w) \in \quantumfields(V)$ be two quantum fields. Then, what is the consecutive action $A(z) B(w)$ ?
        \end{question}
        To answer the question above, pick arbitrary elements $v \in V$ and $\varphi \in V^*$, let:
            $$A(z) := \sum_{i \in \Z} A_i z^{-i}, B(w) := \sum_{j \in \Z} B_j w^{-i}$$
        and then consider the canonical pairing:
            $$\< \varphi , A(z) B(w) v \> := \varphi\left( A_i B_j v \right) z^{-i} w^{-j} \in k(\!(z)\!)(\!(w)\!)$$
        If we were to interpret $[A(z), B(w)] = 0$ as occuring when:
            $$\varphi\left( A_i B_j v \right)z^{-i} w^{-j} = \varphi\left( B_i A_j v \right)w^{-i} z^{-j}$$
        then we would have that:
            $$\< \varphi , A(z) B(w) v \>, \< \varphi , B(w) A(z) v \> \in k(\!(z)\!)(\!(w)\!) \cap k(\!(w)\!)(\!(z)\!) = k[\![z, w]\!][z^{-1}, w^{-1}]$$
        The vector space $k[\![z, w]\!][z^{-1}, w^{-1}]$, however, is strictly smaller than $k(\!(z, w)\!)$, which suggests to us that the interpretation above is slightly too restrictive. Ideally we would like:
            $$\< \varphi , A(z) B(w) v \>, \< \varphi , B(w) A(z) v \> \in k(\!(z, w)\!)$$
        generally, and this is because the consecutive actions of two quantum fields $A, B$ on the same wave function/particle $v \in V$ ought to be supported on the diagonal subscheme of $\Spec k[\![z, w]\!]$, and note that the function field of $\Spec k[\![z, w]\!]$ is nothing but $k(\!(z, w)\!)$.
        
        We are thus led to the following definition:
        \begin{definition}[Locality] \label{def: locality}
            Let $V$ be a vector space. A pair of quantum fields $A(z), B(w) \in \quantumfields(V)$ is said to be \textbf{local} if the orders of the poles of $[A(z), B(w)] := A(z)B(w) - B(z)A(w)$ are uniformly bounded by some $N_{A, B} \in \Z_{\geq 0}$ (depending on $A, B$), i.e. there exists some $N_{A, B} \in \Z_{\geq 0}$ such that:
                $$N \geq N_{A, B} \implies (z - w)^N [A(z), B(w)] := A(z)B(w) - B(z)A(w) = 0$$
            with the expression being understood to be in $\End_k(V)[\![z^{\pm 1}, w^{\pm 1}]\!]$.
        \end{definition}